\textbf{COLUNI 2024 -- Questão 28}\par\medskip Leia o texto a seguir.

A alimentação para corredores é o fator-chave para que o praticante, iniciante ou não, obtenha bons resultados e evite lesões. A corrida é um esporte que exige muita energia na sua prática. Por isso, é necessário suprir os estoques do organismo antes de sair para correr e repor o que foi perdido, logo após o treino.

Na alimentação pré-treino (ou pré-corrida, se preferir) é importante que o corredor procure por alimentos energéticos, como forma de ofertar combustível para o organismo suportar todo o percurso. A banana, a aveia, os cereais integrais, o pão integral e as raízes, como a batata doce ou aipim, são boas opções para serem ingeridas antes de encarar uma corrida.

Vale ressaltar que existem estratégias diferentes que podem ser adotadas por determinados períodos com o objetivo de potencializar e otimizar a utilização de gorduras durante o exercício e, assim, diminuir o consumo do glicogênio (reserva energética do músculo). Essas estratégias, geralmente, utilizam da redução de carboidratos e maior quantidade de gorduras boas antes do treino, como, por exemplo, abacate, oleaginosas, açaí (sem xarope) e até mesmo chocolate amargo.

Na alimentação pós-treino (ou pós-corrida, se preferir) o corredor deve pensar que precisa recuperar o que foi perdido e ainda dar nutrientes para que o organismo consiga reconstruir o que foi utilizado em demasia. Dessa forma, alimentos ricos em proteínas, como barras de proteínas, ovos cozidos, atum, sardinha e carne ou frango, são ótimas fontes alimentares pós-corrida. Para quem não tem alergia ou sensibilidade ao leite, por exemplo, pode utilizá-lo, bem como os seus derivados ou suplementos proteicos.

Disponível em: https://www.oceandrop.com.br/blog/6-dicas-para-corredores-iniciantes. Acesso em: 28 Ago. 2023. (Adaptado)

Com base no texto acima, responda às questões 27 a 30.

O sistema cardiovascular é muito importante para o transporte de nutrientes e $O_2$ pelo corpo de todos os indivíduos. Em atletas, o corpo é mais bem adaptado a esse tipo de esforço físico e capaz de otimizar o trabalho do sistema cardiovascular, permitindo um bom desempenho na atividade esportiva.

\par

Sobre o sistema cardiovascular, marque a afirmativa INCORRETA:\begin{enumerate}
\item O sangue presente nas veias pulmonares é pobre em $O_2$ e rico em $CO_2$.
\item O ventrículo esquerdo do coração bombeia sangue rico em $O_2$ para a grande circulação.
\item As artérias pulmonares transportam sangue pobre em $O_2$ do coração aos pulmões.
\item O sangue pobre em $O_2$ retorna ao coração pelas veias cavas, na grande circulação.\end{enumerate}
